\section{Conceptual Framework}\label{sec:literature}

I organize the mechanism discussion around two frictions that can both distort
household inflation expectations: diagnostic overreaction and information
rigidity. Diagnostic expectations follow representativeness. Agents overweight
states that become vivid in recent news and underweight base rates
\citep{Kahneman1972, Gennaioli2010}. In macro applications, this channel
implies revisions that overshoot and then partially reverse
\citep{Bordalo2018, Bordalo2020b}. Information-rigidity models impose a
different friction. Households update slowly because acquiring and processing
information is costly \citep{Sims2003, Coibion2015}. This channel implies
gradual convergence to the full-information benchmark, not immediate overshoot.

This paper estimates both channels jointly and uses sign and dynamic
restrictions to separate them. The question is not whether one channel exists
in isolation. The question is which channel dominates forecast-error variation
in a setting with high food-price salience and relatively opaque policy
communication. China is informative on this margin. Pork-price shocks,
including African Swine Fever episodes, make inflation highly visible. At the
same time, institutional communication frictions can sustain slow updating
\citep{Chen2022}. These two features can amplify both channels, but they lead
to different policy prescriptions.

\subsection{Diagnostic vs rigidity predictions}\label{sec:diagnostic_vs_rigidity}

Define the one-step-ahead forecast error and forecast revision as
\begin{equation*}
  FE_{t+1} \equiv \pi_{t+1} - \mu_t, \qquad FR_t \equiv \mu_t - \mu_{t-1},
\end{equation*}
where $\pi_{t+1}$ is realized inflation and $\mu_t$ is the expectation formed
at time $t$. I estimate the moments in a nested specification:
\begin{equation*}
  FE_{t+1} = \alpha + \beta FR_t + \lambda FE_t + X_t'\Gamma + u_{t+1}.
\end{equation*}
The coefficient $\beta$ captures revision overshoot. The coefficient $\lambda$
captures persistence in inherited forecast gaps.

Under diagnostic expectations, upward revisions are too large relative to
subsequent realizations. A positive revision is therefore followed by a
negative future error. The key restriction is
\begin{equation}\label{eq:diagnostic_prediction}
  \beta < 0 \;\Longleftrightarrow\; \E[FE_{t+1}\mid FR_t>0] < 0.
\end{equation}
This is the core overreaction moment in \citet{Bordalo2020b}. In dynamic
terms, it implies an overshoot-and-correction profile: revisions move quickly,
then errors partially reverse.

Under information rigidity, households close information gaps slowly. If a
household underpredicts inflation in period $t$, that error tends to persist
into $t+1$ because updating is incomplete. The key restriction is
\begin{equation}\label{eq:rigidity_prediction}
  \lambda > 0 \;\Longleftrightarrow\; \E[FE_{t+1}\mid FE_t>0] > 0.
\end{equation}
This is the standard sluggish-updating moment in \citet{Coibion2015}. Dynamic
adjustment is monotone rather than reversal-based: errors decay gradually while
keeping the same sign across nearby horizons.

The two channels can coexist. A household may underreact to low-salience
background information and overreact when a salient shock dominates attention.
In that case, $\beta<0$ and $\lambda>0$ can both hold in the same sample. The
framework therefore treats mechanism identification as a quantitative problem:
how much each channel contributes, and when. Joint estimation is essential
because separate regressions can misattribute one channel to the other.

The framework also gives rejection criteria. If estimates repeatedly show
$\beta\ge 0$, diagnostic overreaction is not first order in this environment.
If estimates repeatedly show $\lambda\le 0$, slow updating is not first order.
If both restrictions fail, neither benchmark mechanism is sufficient. If both
hold, heterogeneity and dynamics determine which policy lever is more valuable.

\subsection{Chinese institutional context}\label{sec:china_context}

China provides conditions under which both channels can be strong. On the
salience side, food inflation is visible in daily shopping, and pork prices are
the primary amplifier of household inflation narratives. African Swine Fever
episodes created abrupt supply losses and large price spikes. When households
observe these spikes repeatedly, recent food news becomes a dominant signal for
aggregate inflation. This pushes behavior toward representativeness and raises
the probability of diagnostic overreaction.

On the rigidity side, households face a communication environment with less
explicit forward guidance and less transparent policy mapping than benchmark
inflation-targeting regimes. \citet{Chen2022} document institutional features
consistent with this characterization, including shifts in operational focus
and limited precision of policy signaling. When households cannot map policy
messages into a clear inflation path, they update with delay. That delay raises
forecast-error persistence, the signature captured by $\lambda>0$.

These two mechanisms are not competing descriptions of different countries.
They can both be active within China at the same time. Food-price salience can
increase overshoot in revision moments, while communication opacity can sustain
persistent lagged errors. The empirical task is to identify which mechanism
accounts for more variation, in which states, and at which horizons.

Policy implications differ by mechanism. If diagnostic overreaction dominates,
effective communication should counter salience cascades, correct exaggerated
food-price narratives quickly, and anchor expectations around medium-run
inflation fundamentals. If rigidity dominates, policy should improve frequency,
clarity, and predictability of communication to reduce the stock of uninformed
households. Both channels favor better communication, but the optimal design is
different: framing against overreaction versus disclosure to close information
gaps.

\subsection{Testable predictions summary}\label{sec:testable_predictions}

The design evaluates four empirical moments that separate the models.
First, the revision-error sign. Diagnostic models predict $\beta<0$: larger
upward revisions should precede more negative forecast errors
\citep{Bordalo2020b}. Pure rigidity does not require this overshoot moment and
typically implies a weak or nonnegative revision slope once lagged errors are
included.

Second, the lagged-error sign. Rigidity predicts $\lambda>0$: positive errors
today should carry into next-period errors \citep{Coibion2015}. Diagnostic
models allow weak or negative lagged-error dependence when correction follows
overshoot. A positive $\lambda$ supports sluggish updating; a nonpositive
$\lambda$ weakens that channel.

Third, cross-sectional media patterns. Diagnostic models predict stronger
overreaction where inflation news is more salient and cognitively dominant.
Rigidity models predict faster updating where information access and policy
signal quality are better, implying weaker persistence in those groups.
Cross-sectional media variation therefore helps separate salience amplification
from information-delay amplification.

Fourth, dynamic patterns. Diagnostic dynamics are front-loaded: revision shocks
move errors quickly, then attenuate or reverse as overshoot is unwound.
Rigidity dynamics are gradual: errors decay slowly with same-sign persistence
across adjacent horizons. Dynamic profiles therefore provide a second
discriminating margin beyond static sign restrictions.

\subsubsection{CUT heterogeneity}\label{subsub:cut_heterogeneity}

I summarize heterogeneity with a CUT decomposition:
cross-sectional media conditions (C), uncertainty states (U), and temporal
distance from salient shocks (T). Each mechanism implies a distinct CUT
signature.

For C, diagnostic models predict more negative $\beta$ where media salience is
higher. Rigidity models predict lower $\lambda$ where households face better
information access and clearer policy communication.

For U, diagnostic intensity should rise when uncertainty is high because agents
lean more on vivid signals when priors are noisy \citep{Bordalo2018}. Rigidity
also predicts stronger persistence if uncertainty slows processing. Joint
estimation identifies whether uncertainty mainly steepens overreaction
($\beta$ channel) or persistence ($\lambda$ channel).

For T, diagnostics predict an overshoot-then-correction sequence: strongest
effects close to salient news and weaker or reversed effects at longer horizons.
Rigidity predicts monotone convergence without an overshoot requirement. Event-
time dynamics therefore provide a direct CUT test of mechanism dominance.
